\chapter{Návrh firmwaru}

Firmware představuje softwarovou vrstvu, která propojuje jádro emulátoru s reálným hardwarem
zařízení. Stejně jako samotný emulátor, firmware byl realizován v jazyce C, a to s využitím vývojové
sady pico-sdk.

Hlavním úkolem firmwaru je správa běhu emulátoru, zajištění plynulého video a audio výstupu a
obsluha vstupních herních ovladačů. Dále implementuje podporu souborového systému FAT pro čtení a
zápis na SD kartu, což umožňuje práci se soubory jak samotnému firmwaru, tak emulovanému systému
(prostřednictvím rozhraní FWX).

\begin{listing}[htbp]
\begin{minted}{c}
while (true) {
    /* Přečtení vstupu */
    update_joypads(&self->bus.joypad1, &self->bus.joypad2);
    /* Reset přerušení se nemá stat v okamžik zpracování
     * požadavku o resetování */
    disable_interrupts();
    /* Zpracování požadavku na obnovení výchozího stavu emulátoru */
    if (self->reset_required) {
        BusEvent new_ev = {
            .id = BUS_EVENT_RESET,
        };

        bus_add_event(&self->bus, &new_ev);
        self->reset_required = false;
    }

    enable_interrupts();
    /* Zpracování čekajících událostí na sběrnici */
    bus_update(&self->bus);
    /* Krok emulační smyčky (vykreslení jednoho snímku) */
    run_frame(self);
}
\end{minted}
\caption{Hlavní smyčka programu firmwaru.}
\label{lst:main_loop}
\end{listing}

Ve výpisu \ref{lst:main_loop} je uveden kód hlavní smyčky firmwaru. V každé iteraci mikrokontrolér
přečte aktuální stav ovladačů a provede emulaci jednoho úplného snímku obrazu. Před samotnou emulací
snímku jsou zpracovány veškeré události nashromážděné od předchozího cyklu (například požadavek na
načtení nového programu z SD karty). Běh emulace je primárně synchronizován s tokem audio signálu
(viz podkapitola \ref{pp:audio}), což zajišťuje stabilní rychlost 60 snímků za sekundu.

Stisknutí resetovacího tlačítka na základní desce je obsluhováno asynchronně pomocí
hardwarového přerušení. Toto přerušení pouze nastaví příznak \texttt{reset\_required} na hodnotu
\texttt{true}. Díky tomuto přístupu nemusí hlavní program v každém cyklu aktivně číst stav pinu
(tzv. \textit{polling}), čímž se šetří cenné procesorové instrukce a vymezuje situaci, když
stisknutí nebylo zaregistrováno.

Z rychlostních důvodu, RP2350 je přetaktován (\textit{overclocked}) z původních 150 MHz na 300 MHz.
